\section{Introduction}
\label{sec:introduction}

The rapid spread of scientific claims through social media creates a difficult fact-checking scenario: posts often report, simplify or sensationalize research findings while omitting the original source. In many cases, the cited study is not linked explicitly with a URL or a bibliographic reference but only mentioned indirectly through fragments such as the reported finding, the venue, the institution or the authors. As a consequence, verifying such claims requires first identifying the scientific publication that is being implicitly referenced \cite{10.1007/978-3-032-21321-1_43}.

This problem is the focus of Task 1 of the CLEF 2026 CheckThat! Lab, \emph{Source Retrieval for Scientific Web Claims}: given a social media post containing a scientific claim and an implicit reference to a paper, the goal is to retrieve the correct publication from a fixed collection of 10,000 candidate papers written in English \cite{clef-checkthat:2026:task1}.

The task is challenging for several reasons. First, the linguistic gap between social media and scientific writing is substantial: posts are short, informal and often contain paraphrases, slogans, hashtags or incomplete statements whereas scientific papers use more precise and technical language. Second, the reference to the source paper is implicit rather than explicit so the retrieval system must infer the connection from partial semantic evidence rather than from direct citations. Third, relevant clues may appear in different forms, such as mentions of a scientific result, a medical outcome, a research institution or a publication venue, requiring the system to combine lexical and semantic evidence effectively.

An additional challenge is multilinguality: the official query sets include English, German and French posts while the candidate pool remains shared across languages.

Team RETRIX approaches this problem with a four-stage sequential pipeline.
\begin{enumerate}
    \item In a preliminary stage we translate to English all the queries written in French and German using the model utter-project/EuroLLM-9B-Instruct-2512 \cite{utter_project_eurollm_9b_instruct_2512}. This allows us to handle all queries in the same way, independently of the original language and consequently of the semantic and structural differences.

    \item Raw tweet-style queries are cleaned of emojis and symbols like "\#" and "@" and then enriched via pseudo-relevance feedback: BAAI/bge-large-en-v1.5 \cite{bge-large-en-v1.5} encodes both queries and corpus documents into a shared embedding space, retrieves the top-5 semantically similar papers per query and extracts the most frequent terms from those papers as expansion tokens.

    \item The expanded queries are fed into BAAI/bge-m3 \cite{bge-m3} operating in hybrid mode, which combines dense inner-product search via FAISS \cite{faiss_documentation} with sparse lexical matching and ColBERT multi-vector scoring to produce an initial ranked list of candidates.

    \item Finally, the top candidates are re-ranked by a neural cross-encoder, in our case nvidia/llama-nemotron-rerank-1b-v2 \cite{llama-nemotron-rerank-1b-v2}, fine-tuned on the English training dataset of the CLEF CheckThat! Task 1 2026 and all the queries' datasets of the 2025 CLEF CheckThat! SubTask 4b.
\end{enumerate}

The paper is organized as follows: Section~\ref{sec:related} discusses the most relevant research directions related to this task; Section~\ref{sec:methodology} describes in more detail our approach; Section~\ref{sec:setup} explains our experimental setup; Section~\ref{sec:results} discusses our main findings; finally, Section~\ref{sec:conclusion} draws some conclusions and outlooks for future work. 